\subsection{The full ramified local metabelian section space}

\begin{statement}{E.66}{Proved}{Every ramified local metabelian section is cuspidal, with its unique full lift}
Retain the curve, field, good place, integral tangents, actual paths
and tame generators of E.61. In particular
$K=\mathbb Q(\zeta_3)$, $v\nmid3$, $q=|k_v|$, and
$\Pi=\pi_1(U_{\overline K})^{(3)}$ is free pro-three of rank ten.
All derived subgroups below are closed. Put
\[
 M=\Pi/\Pi'',\qquad V=\Pi^{\rm ab},\qquad
 A=\Pi'/\Pi'',\qquad v_i=[d_i]\in V.
 \tag{E66.1}
\]
These are actual continuous quotients with the genuine unramified
action fixed in E.61. The space of local $M$-section classes with
nontrivial geometric inertia is exactly
\[
 \coprod_{i=1}^{9}
   (\mathbb Z_3\setminus\{0\})\times
                \mathbb Z_3/(q-1)\mathbb Z_3.
 \tag{E66.2}
\]
The representative for $(i,a,\bar b)$ is the image in $M$ of
the actual full tangential section
\[
            t\longmapsto d_i^a t,\qquad
            F\longmapsto d_i^bF_i.
 \tag{E66.3}
\]
Every equivalence in (E66.2) is conjugation by $d_i^k$, sending
$b$ to $b+(1-q)k$. The cusp and the nonzero exponent are unchanged.
The parametrization is a homeomorphism for the quotient topology.

The actual map
\[
 \operatorname{Sect}(\Pi\rtimes G_{K_v})/\Pi
       \longrightarrow
 \operatorname{Sect}(M\rtimes G_{K_v})/M
 \tag{E66.4}
\]
is a continuous surjection. Its restriction over the ramified locus
is a homeomorphism; in particular every ramified metabelian class
has exactly one full section class above it. Its unramified locus
also has full lifts, by lifting its arbitrary Frobenius coordinate.
The further classification of those unramified conjugacy fibers is
retained in the specified remaining local and global maps of E.23.

For the actual quotient $M\twoheadrightarrow R$ at
$K_v=\mathbb Q_2(\zeta_3)$, the ramified image therefore has
precisely the cusp representatives of E.61. The displayed
$R$-section of E61.5 lies outside that image. All nine cusp
inertia lines, its twenty-seven finite comparisons and the
detecting degree-$19683$ cover retain their literal maps and
certificates. The theorem settles the complete ramified local
metabelian source at every specified good place, with no
theta-value normalization required.
The manuscript anchor is \S51.16, folios 733v--736v,
scans172--178, formulas(371)--(379), Part E(c),(d).
\end{statement}

\begin{proof}
\emph{The arithmetic relation and its finite covers.}
E.61 identifies the actual good-reduction action with the special
fibre action and proves its pro-three factorization. A local
section with geometric kernel $M$ is a pair
$(\bar\nu,\bar z)\in M^2$ satisfying
\[
        \bar zF(\bar\nu)\bar z^{-1}=\bar\nu^q.
 \tag{E66.5}
\]
Its abelian inertia $w\in V$ satisfies $F(w)=qw$.
E61.14 proves that $w$ belongs to one of the nine integral
cusp lines. We must determine the entire inertia element,
including the case $w=0$.

Choose any lift $z\in\Pi$ of $\bar z$, and write $\psi=zF$.
For an $F$-stable open lattice $W_n\subset V$, let
$H_n$ be its inverse image in $\Pi$. The group $H_n$ contains
$\Pi'$, hence $\Pi''\subset[H_n,H_n]$. It is $\psi$-stable.
As in E61.7--E61.11, the actual open arithmetic subgroup
$H_n\rtimes\langle\psi\rangle$ defines a finite etale cover
$U_n\to U_{k_v}$ over $k_v$, with smooth compactification $C_n$.
All cover maps used below are the inclusion maps between these
subgroups, with this same Frobenius lift $\psi$.

When $w\in W_n$, the image $\nu_n\in H_n^{\rm ab}$ of a lift
of $\bar\nu$ is independent of that lift. Relation (E66.5)
gives the exact equality
\[
                  \psi(\nu_n)=q\nu_n
                    \quad\hbox{in }H_n^{\rm ab}.
 \tag{E66.6}
\]
The continuous, Frobenius-equivariant filling sequence is
\[
 0\longrightarrow B_n\longrightarrow H_n^{\rm ab}
       \longrightarrow T_3\operatorname{Jac}(C_n)
       \longrightarrow0,
 \tag{E66.7}
\]
where $B_n$ is the free cusp lattice modulo its single diagonal
relation. The integral exactness follows from the punctured-surface
presentation, with its oriented cusp generators; its specialization
and Jacobian comparison are those proved in E.61 from
\cite[XIII, 2.9--2.12]{E_SGA1} and
\cite[III, Proposition9.1]{E_MilneAV2008}.
The finite-field Frobenius eigenvalues on the final Tate module
have complex absolute value $\sqrt q$
\cite[II, Theorem1.1(b)]{E_MilneAV2008}.
It follows from (E66.6) that
\[
                         \nu_n\in B_n .
 \tag{E66.8}
\]
This application uses the actual finite-field cover $C_n$;
no weight statement about an arbitrary local-field Jacobian is
substituted for that comparison.

\emph{Two exact inverse-limit comparisons.}
We record the maps needed to use all the finite covers at once.
For a decreasing sequence of closed subgroups $H_n$ of the
compact group $\Pi$, put $H_\infty=\bigcap_nH_n$.
The natural inclusion maps give an isomorphism
\[
                 H_\infty^{\rm ab}
                       \xrightarrow{\ \sim\ }
                       \varprojlim_n H_n^{\rm ab}.
 \tag{E66.9}
\]
Here abelianization is the maximal abelian pro-three quotient.
To prove this, take a continuous character of $H_\infty$ with
finite three-primary image. Its kernel contains
$H_\infty\cap U$ for some open normal subgroup $U$ of $\Pi$.
The decreasing images of the $H_n$ in the finite group $\Pi/U$
eventually equal the image of $H_\infty$: each element in their
intersection lifts to the intersection by compactness.
The character therefore extends to all sufficiently deep $H_n$.
If such a character restricts to zero on $H_\infty$, choose $U$
inside its kernel on the relevant $H_n$ and use the same finite
image argument to see that it is zero on every sufficiently
deep $H_m$. Thus restriction identifies the direct limit of the
finite-character groups with the character group of
$H_\infty$. Taking their compact Pontryagin duals proves
(E66.9), with exactly the displayed inclusion maps.

We also take inverse limits of compact short exact sequences
\[
       0\longrightarrow\Delta_n\longrightarrow P_n
                \longrightarrow B_n\longrightarrow0 .
 \tag{E66.10}
\]
Exactness at the right can be proved directly. For a compatible
sequence in the $B_n$, the set of lifts satisfying any finite
collection of compatibility equations is nonempty: lift its
coordinate at the largest index and map that lift down.
These are closed subsets of the compact product of the $P_n$.
Their finite intersection property supplies a compatible lift.
The kernel is the inverse limit of the displayed $\Delta_n$.
This proves exactness here without a surjectivity assumption
on the transition maps of the diagonal lattices.

\emph{Vanishing of metabelian inertia with zero abelian image.}
Suppose first that $w=0$, and set $W_n=3^nV$.
Then $H_\infty=\Pi'$. The geometric deck group of the cover is
$Q_n=V/3^nV$. Write $\bar I_{i,n}$ for the subgroup generated
by $v_i$ in this finite group. Since every $v_i$ is primitive,
it has order $3^n$. The exact cusp modules in (E66.10) are
\[
 P_n=\bigoplus_{i=1}^{9}\mathbb Z_3[Q_n/\bar I_{i,n}],
 \qquad
 \Delta_n=\mathbb Z_3\sum_{i,g}e_{i,g}.
 \tag{E66.11}
\]
The basis vector $e_{i,g}$ is the actual oriented inertia
class of $g d_i^{3^n}g^{-1}$.
Under $H_{n+1}\hookrightarrow H_n$, every such cusp class
maps to three times the corresponding class below.
Consequently the transition on each summand of $P_n$ is
three times the trace along its finite coset map.
Every fixed coordinate of a compatible inverse sequence is
therefore divisible by every power of three; thus
$\varprojlim P_n=0$. Exactness of (E66.10) gives
$\varprojlim B_n=0$ as well.
By (E66.8)--(E66.9), the class of $\bar\nu$ in
$\Pi'/\Pi''$ is zero. Hence $\bar\nu=1$ in $M$.

\emph{The surviving cusp module for nonzero abelian inertia.}
Now let $w=a v_i\ne0$, put $r=v_3(a)$ and
$u=a/3^r\in\mathbb Z_3^\times$, and use
\[
 W_n=\mathbb Z_3w+3^nV\quad(n\ge r),\qquad
 Y=V/\mathbb Z_3v_i .
 \tag{E66.12}
\]
We have $H_\infty=\Pi'\langle d_i^{3^r}\rangle$.
Any two distinct cusp vectors occur together in a basis of
$V$, by the punctured-surface basis used in E.61.
In $V/W_n$, cusp $i$ has order $3^r$, whereas every other
cusp has order $3^n$. Thus its own boundary summand has
transition equal to the trace
\[
 \mathbb Z_3[Y/3^{n+1}Y]\longrightarrow
                         \mathbb Z_3[Y/3^nY],
 \tag{E66.13}
\]
and each of the other eight transitions is three times a trace.
Their inverse limits are zero by the preceding divisibility
argument. The covering degree between consecutive levels is
$3^9$, so the transition on $\Delta_n$ is multiplication
by $3^9$ and its inverse limit is zero. The exact compact
comparison (E66.10) therefore identifies the full boundary limit
as
\[
             \varprojlim_nB_n
                  \simeq\mathbb Z_3[[Y]] .
 \tag{E66.14}
\]
The right side is the module of integral three-adic measures
on the compact additive group $Y$. This identification sends
the inertia class of $g d_i^{3^r}g^{-1}$ to the Dirac measure
at the image of $g$ in $Y$. Thus it specifies the actual
geometric origin of every basis coordinate.

The Frobenius action on this module is $q$ times the
permutation of measures induced by the affine map
\[
       T:Y\longrightarrow Y,
       \qquad \bar g\longmapsto
          F(\bar g)+\overline{[z]-[z_i]} .
 \tag{E66.15}
\]
Indeed $F_i=z_iF$ sends $d_i$ to $d_i^q$; conjugating the
displayed cusp word by $zF$ gives precisely the label and
factor in (E66.15). This verifies both the translation term
and the orientation factor, without an arbitrary choice of
a nonbase cusp Frobenius origin.

\emph{Integral invariant measures for this actual affine action.}
The eigenvalues of $F$ on $Y\otimes\mathbb Q_3$ are $q$ on
the remaining seven-dimensional cusp space and the two
finite-field elliptic eigenvalues of complex absolute value
$\sqrt q$ on the compact quotient. In particular $F^m-I$
is invertible over $\mathbb Q_3$ for every positive integer $m$.
The affine map (E66.15) has one fixed point in
$Y\otimes\mathbb Q_3$, obtained by applying $(I-F)^{-1}$
to its translation. Every affine periodic point equals
this fixed point: the equation for period $m$ has a unique
solution and that same fixed point is a solution.

Moreover the closure of the permutation group generated by
$T$ is pro-three. The linear action is pro-three by E.61,
and translations form the pro-three group $Y$.
Thus every orbit on every finite quotient $Y/3^nY$ has
three-power length.

Let $\mu\in\mathbb Z_3[[Y]]$ be $T$-invariant.
Take an orbit $O$ on a fixed finite quotient whose full
inverse image does not contain the rational fixed point.
For each integer $b$, all the orbits above $O$ on a
sufficiently deep quotient have length at least
$3^b|O|$. To see this, the closed sets on that inverse
image defined by $T^{3^c}x\equiv x\pmod{3^nY}$ decrease,
as $n$ grows, to its fixed-point set for $T^{3^c}$.
This set is empty. Compactness gives a uniform depth
where they are empty, for each prescribed $c$.

At a deep enough level, the contribution of each refined
orbit to a fixed coarse coset is its orbit-length ratio
times its common integral coefficient. It is divisible
by $3^b$. Trace compatibility therefore forces the
coefficient of every coset of $O$ to be divisible by every
power of three, hence to vanish. All clopen sets avoiding
the fixed point have measure zero.
It follows that
\[
 (\mathbb Z_3[[Y]])^T=
 \begin{cases}
  \mathbb Z_3\delta_{g_0},&g_0\in Y,\\
  0,&g_0\notin Y,
 \end{cases}
 \tag{E66.16}
\]
where $g_0=(I-F)^{-1}\overline{[z]-[z_i]}$ is the stated
rational fixed point. When it lies outside $Y$, every
coarse coset avoids it, so the same argument kills all
coefficients.

Apply this calculation to the compatible boundary class
of $\bar\nu$ supplied by (E66.8). Relation (E66.6), together
with the factor $q$ in (E66.15), says that its measure is
$T$-invariant. Its total mass is $u$: under the inclusion
in $V$, each selected cusp generator maps to $3^rv_i$,
whereas $\bar\nu$ maps to $a v_i$. Thus $g_0\in Y$ and
the measure is exactly $u\delta_{g_0}$.
Choose $g\in\Pi$ with this image in $Y$. The elements
$\bar\nu$ and $g d_i^a g^{-1}$ have the same image in
$H_\infty^{\rm ab}$, by (E66.9) and (E66.14).

\emph{The exact conjugating correction.}
In $M$ the subgroup $H_\infty/\Pi''$ is the split extension
$A\rtimes\langle d_i^{3^r}\rangle$. Its derived group is
the closed image
\[
 [H_\infty,H_\infty]/\Pi''
           =(d_i^{3^r}-1)A=(d_i^a-1)A .
 \tag{E66.17}
\]
The first equality follows by multiplying commutators in
this split extension, with $A$ abelian. Its displayed image
is closed because $A$ is compact. For the second equality,
the formal power series $((1+T)^u-1)/T$ is a unit in
$\mathbb Z_3[[T]]$, since its constant coefficient is
$u\in\mathbb Z_3^\times$. Evaluate it at the action of
$d_i^{3^r}-1$ on $A$.
Conjugation of $g d_i^a g^{-1}$ by elements of $A$ changes
its $A$ coordinate by exactly $(1-d_i^a)A$.
The agreement in $H_\infty^{\rm ab}$ and (E66.17)
therefore supply an actual $A$-conjugator carrying it to
$\bar\nu$. We have proved that every nonzero metabelian
inertia is conjugate to one of the actual $d_i^a$.

\emph{Centralizers and all equivalences.}
We verify the required centralizer in the metabelian group,
including nonunit exponents. Choose a free pro-three basis
$d_i=b_1,b_2,\ldots,b_{10}$ of $\Pi$, put
$\Lambda=\mathbb Z_3[[V]]$, and write $X_j$ for its basis
elements. The continuous Magnus map is
\[
 M\longrightarrow V\ltimes\Lambda^{10},\qquad
 b_j\longmapsto(X_j,e_j),
 \qquad (X,c)(Y,d)=(XY,c+Xd).
 \tag{E66.18}
\]
It is injective. Here is the precise comparison proving
injectivity. For the open subgroups above $3^nV$, the
finite covering graph of the ten-petal rose identifies
$H_n^{\rm ab}$ with the kernel of its integral edge-boundary
map
$\mathbb Z_3[V/3^nV]^{10}\to\mathbb Z_3[V/3^nV]$.
The maps between these graphs are the actual edge-trace
maps. Taking their compact inverse limit and using
(E66.9) embeds $A$ in $\Lambda^{10}$. The Magnus map
restricts to this embedding on $A$, while its first
coordinate is the identity on $V$; hence its full kernel
is trivial. Its path-boundary identity is
\[
                 \sum_j(X_j-1)c_j=X-1 .
 \tag{E66.19}
\]

For an element $(X,c)$ commuting with $d_i^a$, multiplication
in (E66.18) gives
\[
 (X_1^a-1)c=(X-1)
       \frac{X_1^a-1}{X_1-1}e_1 .
 \tag{E66.20}
\]
The ring $\Lambda\simeq\mathbb Z_3[[T_1,\ldots,T_{10}]]$
is a domain, and $a\ne0$ makes the first factor nonzero.
Thus $c_j=0$ for $j>1$. Identity (E66.19), reduced at
$X_1=1$, forces $X=X_1^b$ for some $b\in\mathbb Z_3$.
It then forces $c_1=(X_1^b-1)/(X_1-1)$.
Injectivity of (E66.18) proves
\[
                         C_M(d_i^a)=\langle d_i\rangle.
 \tag{E66.21}
\]
This calculation never divides by the possibly nonunit
integer $a$.

Conjugate the section so that $\bar\nu=d_i^a$.
Both its Frobenius lift and $F_i$ conjugate $d_i^a$ to
$d_i^{qa}$. Their ratio centralizes this nonzero power;
(E66.21) makes the ratio $d_i^b$. Thus the section is
exactly (E66.3). Any geometric conjugation between two
such representatives preserves the cusp and $a$, since
their abelian cusp vectors do. It centralizes that
inertia, hence lies in $\langle d_i\rangle$, and changes
$b$ by precisely $(1-q)k$. This proves all the classes
and equivalences in (E66.2).
The same compact-annulus argument as in E.61 proves the
homeomorphism: on each clopen locus $v_3(a)=r$ the
parametrization is a continuous bijection from a compact
space to a Hausdorff geometric-conjugacy quotient.

Finally the full ramified section classification E.61
has these identical parameters and representatives, so
(E66.4) is a homeomorphism on that locus and each fiber
there is a single class. A section with $\bar\nu=1$ has
an arbitrary Frobenius coordinate $\bar z\in M$.
Choose any lift $z\in\Pi$. The pair $(1,z)$ satisfies
the full tame relation and defines a continuous full
section by E61.8. This proves surjectivity everywhere.
A full ramified section has nonzero abelian cusp inertia,
so its image in $M$ remains ramified; the unramified
lifting assertion retains precisely the unramified
conjugacy fibers stated in the theorem.
\end{proof}
